\scanpage{18}
% L1–4
conséquence de notre hypothèse), et \deleted{\unsure{que}} son rang sera nécessairement $p$. Ceci dit il faut exprimer que les $D^i$ $(0\leq i\leq p-1)$ forment une \uline{base} du $\unsure{K|C|k}$-module $\operatorname{Hom}_A(C,C)$ [où on pose $(\lambda u)(\unsure{\mu|\rho|y})=\lambda(u(\mu))$].

% L5
\deleted{Soit p.ex.}

% L6–8
Supposons p.ex. que $C$ admette une $p$-base \inserted{sur $A$} réduite à un élément $u$, donc \unsure{on a|que}
\[
C\simeq A[u]/(u^p-a),\qquad\text{où }a\in A.
\]
\deleted{Posons} \deleted{\unsure{Donc|$Du$}} Alors une $A$-dérivation de $C$ est connue quand on connaît sa valeur sur $u$, qui peut-être choisie arbitrairement.

% L12–14
Définissons alors la $A$-dérivation $D$ dans $C$ par
\[
\begin{aligned}
&Du=1\\
\text{d'où}\qquad &Du^n=nu^{n-1}
\end{aligned}
\]
% L15–19
On a donc $D^p=0$, donc $C$ est un $\deleted{\unsure{D}}\unsure{U|\mathcal U}$-module. De plus, \deleted{la matrice des} identifions, grâce à la base $(u^i)_{0\leq i\leq p-1}$, \deleted{\unsure{le|et}} $\operatorname{Hom}(C,C)$ à $C^p$, par $v\mathrel{\unsure{\text{\textemdash}|\longrightarrow|\mapsto}}(v(u^i))_{\unsure{0\leq i\leq p-1}}$ [c'est un isom. de $C$-modules]. Alors
% L20–21
la matrice du système des $D^i$ $(0\leq i\leq p-1)$ est
\[
\left(\begin{array}{ccccc||l}
1&&u&u^2&u^3&u^{p-1}\\
0&&1&2u&3u^2&(p-1)u^{p-2}\\
0&&0&2.1&3.2.u&(p-1)(p-2)u^{p-3}\\
0&&0&0&3.2.1.&(p-1)(p-2)(p-3)u^{p-4}\\ \cline{2-6}
0&0&0&0&0&(p-1)\cdots1.
\end{array}\right)
\]
% archive sheet number: 17
